\chapter{Literature Review}

\section{Related Work}
\TemplateTodo{Summarize the most relevant related studies. Explain their objectives, methods, data, results, and limitations. End by positioning your study against those works.}

A literature review is not merely a list of article summaries. Organize it thematically so readers understand the development of knowledge, remaining gaps, and why this study is needed.

\begin{table}[H]
\centering
\caption{Example format for summarizing related work}
\begin{tabularx}{\linewidth}{@{}p{0.18\linewidth}X X@{}}
\toprule
Author & Focus and Method & Relevance to This \DocTypeName \\
\midrule
Author name (year) & Short summary of method, data, and main results. & Explain lessons, gaps, or distinctions for your research. \\
\bottomrule
\end{tabularx}
\end{table}

\section{Theoretical Foundation}
\TemplateTodo{Explain concepts, theories, algorithms, standards, or frameworks that are actually used in the study.}

\subsection{First Main Concept}
Write the definition, working principle, and relevance of the concept to the study. Use citations for definitions or theories that are not your own opinion.

\subsection{Second Main Concept}
Add other concepts as needed. If using formulas, explain the meaning of each symbol and why the formula is used.

\section{Conceptual Framework}
\TemplateTodo{Explain the relationship between the problem, theory, method, data, and research outputs. A conceptual diagram can help this section.}

\begin{figure}[H]
\centering
\includegraphics[width=0.75\linewidth]{example-figure.pdf}
\caption{Example figure usage in the template}
\label{fig:example-figure}
\end{figure}

An example of figure usage is shown in Figure~\ref{fig:example-figure}. Mathematical notation can be used to explain the relationship between layers in Figure~\ref{fig:example-figure}. Let the input vector be denoted as
\begin{equation}
\mathbf{x} =
\begin{bmatrix}
x_1 & x_2 & \cdots & x_n
\end{bmatrix}^{\mathsf{T}}.
\end{equation}
For the $l$-th hidden layer, neuron activations are computed by
\begin{equation}
\mathbf{a}^{[l]} = \sigma\left(\mathbf{W}^{[l]}\mathbf{a}^{[l-1]} + \mathbf{b}^{[l]}\right),
\qquad l = 1,2,\ldots,L-1,
\end{equation}
where $\mathbf{a}^{[0]}=\mathbf{x}$, $\mathbf{W}^{[l]}$
is the weight matrix at layer $l$,
$\mathbf{b}^{[l]}$ is the bias vector, and
$\sigma(\cdot)$ is the activation function. The model output is denoted as
\begin{equation}
\hat{\mathbf{y}} = f\left(\mathbf{W}^{[L]}\mathbf{a}^{[L-1]} + \mathbf{b}^{[L]}\right),
\end{equation}
with output components
\begin{equation}
\hat{\mathbf{y}} =
\begin{bmatrix}
\hat{y}_1 & \hat{y}_2 & \cdots & \hat{y}_k
\end{bmatrix}^{\mathsf{T}}.
\end{equation}
